% Options for packages loaded elsewhere
% Options for packages loaded elsewhere
\PassOptionsToPackage{unicode}{hyperref}
\PassOptionsToPackage{hyphens}{url}
\PassOptionsToPackage{dvipsnames,svgnames,x11names}{xcolor}
%
\documentclass[
  british,
]{article}
\usepackage{xcolor}
\usepackage{amsmath,amssymb}
\setcounter{secnumdepth}{5}
\usepackage{iftex}
\ifPDFTeX
  \usepackage[T1]{fontenc}
  \usepackage[utf8]{inputenc}
  \usepackage{textcomp} % provide euro and other symbols
\else % if luatex or xetex
  \usepackage{unicode-math} % this also loads fontspec
  \defaultfontfeatures{Scale=MatchLowercase}
  \defaultfontfeatures[\rmfamily]{Ligatures=TeX,Scale=1}
\fi
\usepackage{lmodern}
\ifPDFTeX\else
  % xetex/luatex font selection
\fi
% Use upquote if available, for straight quotes in verbatim environments
\IfFileExists{upquote.sty}{\usepackage{upquote}}{}
\IfFileExists{microtype.sty}{% use microtype if available
  \usepackage[]{microtype}
  \UseMicrotypeSet[protrusion]{basicmath} % disable protrusion for tt fonts
}{}
\makeatletter
\@ifundefined{KOMAClassName}{% if non-KOMA class
  \IfFileExists{parskip.sty}{%
    \usepackage{parskip}
  }{% else
    \setlength{\parindent}{0pt}
    \setlength{\parskip}{6pt plus 2pt minus 1pt}}
}{% if KOMA class
  \KOMAoptions{parskip=half}}
\makeatother
% Make \paragraph and \subparagraph free-standing
\makeatletter
\ifx\paragraph\undefined\else
  \let\oldparagraph\paragraph
  \renewcommand{\paragraph}{
    \@ifstar
      \xxxParagraphStar
      \xxxParagraphNoStar
  }
  \newcommand{\xxxParagraphStar}[1]{\oldparagraph*{#1}\mbox{}}
  \newcommand{\xxxParagraphNoStar}[1]{\oldparagraph{#1}\mbox{}}
\fi
\ifx\subparagraph\undefined\else
  \let\oldsubparagraph\subparagraph
  \renewcommand{\subparagraph}{
    \@ifstar
      \xxxSubParagraphStar
      \xxxSubParagraphNoStar
  }
  \newcommand{\xxxSubParagraphStar}[1]{\oldsubparagraph*{#1}\mbox{}}
  \newcommand{\xxxSubParagraphNoStar}[1]{\oldsubparagraph{#1}\mbox{}}
\fi
\makeatother


\usepackage{longtable,booktabs,array}
\newcounter{none} % for unnumbered tables
\usepackage{calc} % for calculating minipage widths
% Correct order of tables after \paragraph or \subparagraph
\usepackage{etoolbox}
\makeatletter
\patchcmd\longtable{\par}{\if@noskipsec\mbox{}\fi\par}{}{}
\makeatother
% Allow footnotes in longtable head/foot
\IfFileExists{footnotehyper.sty}{\usepackage{footnotehyper}}{\usepackage{footnote}}
\makesavenoteenv{longtable}
\usepackage{graphicx}
\makeatletter
\newsavebox\pandoc@box
\newcommand*\pandocbounded[1]{% scales image to fit in text height/width
  \sbox\pandoc@box{#1}%
  \Gscale@div\@tempa{\textheight}{\dimexpr\ht\pandoc@box+\dp\pandoc@box\relax}%
  \Gscale@div\@tempb{\linewidth}{\wd\pandoc@box}%
  \ifdim\@tempb\p@<\@tempa\p@\let\@tempa\@tempb\fi% select the smaller of both
  \ifdim\@tempa\p@<\p@\scalebox{\@tempa}{\usebox\pandoc@box}%
  \else\usebox{\pandoc@box}%
  \fi%
}
% Set default figure placement to htbp
\def\fps@figure{htbp}
\makeatother


% definitions for citeproc citations
\NewDocumentCommand\citeproctext{}{}
\NewDocumentCommand\citeproc{mm}{%
  \begingroup\def\citeproctext{#2}\cite{#1}\endgroup}
\makeatletter
 % allow citations to break across lines
 \let\@cite@ofmt\@firstofone
 % avoid brackets around text for \cite:
 \def\@biblabel#1{}
 \def\@cite#1#2{{#1\if@tempswa , #2\fi}}
\makeatother
\newlength{\cslhangindent}
\setlength{\cslhangindent}{1.5em}
\newlength{\csllabelwidth}
\setlength{\csllabelwidth}{3em}
\newenvironment{CSLReferences}[2] % #1 hanging-indent, #2 entry-spacing
 {\begin{list}{}{%
  \setlength{\itemindent}{0pt}
  \setlength{\leftmargin}{0pt}
  \setlength{\parsep}{0pt}
  % turn on hanging indent if param 1 is 1
  \ifodd #1
   \setlength{\leftmargin}{\cslhangindent}
   \setlength{\itemindent}{-1\cslhangindent}
  \fi
  % set entry spacing
  \setlength{\itemsep}{#2\baselineskip}}}
 {\end{list}}
\usepackage{calc}
\newcommand{\CSLBlock}[1]{\hfill\break\parbox[t]{\linewidth}{\strut\ignorespaces#1\strut}}
\newcommand{\CSLLeftMargin}[1]{\parbox[t]{\csllabelwidth}{\strut#1\strut}}
\newcommand{\CSLRightInline}[1]{\parbox[t]{\linewidth - \csllabelwidth}{\strut#1\strut}}
\newcommand{\CSLIndent}[1]{\hspace{\cslhangindent}#1}

\ifLuaTeX
\usepackage[bidi=basic,shorthands=off]{babel}
\else
\usepackage[bidi=default,shorthands=off]{babel}
\fi
\ifLuaTeX
  \usepackage{selnolig} % disable illegal ligatures
\fi


\setlength{\emergencystretch}{3em} % prevent overfull lines

\providecommand{\tightlist}{%
  \setlength{\itemsep}{0pt}\setlength{\parskip}{0pt}}



 


% Get Math Fonts
\usepackage{fontspec}
\usepackage{unicode-math}
\setmathfont{Latin Modern Math}
\setmathfont[range={\mathscr,\mathbfscr,\mathbb}]{XITSMath-Regular}
[    Extension = .otf,
      BoldFont = XITSMath-Bold
]
\makeatletter
\@ifpackageloaded{caption}{}{\usepackage{caption}}
\AtBeginDocument{%
\ifdefined\contentsname
  \renewcommand*\contentsname{Table of contents}
\else
  \newcommand\contentsname{Table of contents}
\fi
\ifdefined\listfigurename
  \renewcommand*\listfigurename{List of Figures}
\else
  \newcommand\listfigurename{List of Figures}
\fi
\ifdefined\listtablename
  \renewcommand*\listtablename{List of Tables}
\else
  \newcommand\listtablename{List of Tables}
\fi
\ifdefined\figurename
  \renewcommand*\figurename{Figure}
\else
  \newcommand\figurename{Figure}
\fi
\ifdefined\tablename
  \renewcommand*\tablename{Table}
\else
  \newcommand\tablename{Table}
\fi
}
\@ifpackageloaded{float}{}{\usepackage{float}}
\floatstyle{ruled}
\@ifundefined{c@chapter}{\newfloat{codelisting}{h}{lop}}{\newfloat{codelisting}{h}{lop}[chapter]}
\floatname{codelisting}{Listing}
\newcommand*\listoflistings{\listof{codelisting}{List of Listings}}
\usepackage{amsthm}
\theoremstyle{plain}
\newtheorem{theorem}{Theorem}[section]
\theoremstyle{remark}
\AtBeginDocument{\renewcommand*{\proofname}{Proof}}
\newtheorem*{remark}{Remark}
\newtheorem*{solution}{Solution}
\newtheorem{refremark}{Remark}[section]
\newtheorem{refsolution}{Solution}[section]
\makeatother
\makeatletter
\makeatother
\makeatletter
\@ifpackageloaded{caption}{}{\usepackage{caption}}
\@ifpackageloaded{subcaption}{}{\usepackage{subcaption}}
\makeatother

\newcommand{\AmenableGroup}{{\Gamma}}
\newcommand{\AmenableGroupElement}{{\gamma}}
\newcommand{\CountingMeasure}[1]{{\left|#1\right|}}
\newcommand{\Folner}[1][\vphantom{}]{{\Phi}_{#1}}
\newcommand{\FurstenbergFolner}[1][{\vphantom{}}]{{\Psi}_{#1}}
\newcommand{\Group}{{\Gamma}}
\newcommand{\GroupAction}[2]{T^{#1}{#2}}
\newcommand{\GroupOperation}[2]{{#1}\cdot{#2}}
\newcommand{\Measure}{{\mu}}
\newcommand{\ShiftSpace}{{X}}

\usepackage{bookmark}
\IfFileExists{xurl.sty}{\usepackage{xurl}}{} % add URL line breaks if available
\urlstyle{same}
% fallback for those not using the hyperref driver hyperxmp:
\makeatletter
\@ifundefined{xmpquote}{\newcommand{\xmpquote}[1]{#1}}{}
\makeatother
\hypersetup{
  pdftitle={Furstenberg's Correspondence Principle},
  pdflang={en-GB},
  colorlinks=true,
  linkcolor={blue},
  filecolor={Maroon},
  citecolor={Blue},
  urlcolor={Blue},
  pdfcreator={LaTeX via pandoc}}


\title{Furstenberg's Correspondence Principle}
\author{}
\date{2026-09-27}
\begin{document}
\maketitle


\begin{theorem}[{cf. Kra et al., 2022, Theorem
2.10}]\protect\hypertarget{thm-FCP}{}\label{thm-FCP}

Let \(\AmenableGroup\) be an amenable group,
\(A\subset \AmenableGroup\), and \(\Folner\) be a Følner sequence in
\(\AmenableGroup\) such that the limit
\[\delta=\lim_{N\rightarrow\infty}\frac{\CountingMeasure{A\cap\Folner[N]}}{\CountingMeasure{\Folner[N]}}\]
exists.

Then there exists an ergodic system
\((\ShiftSpace,\Measure,\AmenableGroup)\) that is acted on by
\(\AmenableGroup\), a clopen set \(E\subset\ShiftSpace\), a Følner
sequence \(\FurstenbergFolner\) in \(\AmenableGroup\), and a point
\(a\in\ShiftSpace\) that is generic with \(\Measure\) with respect to
\(\FurstenbergFolner\) such that \(\Measure(E)\geq\delta\) and
\[A=\{\AmenableGroupElement\in \AmenableGroup:\GroupAction{\AmenableGroupElement}{a}\in E\}. \]

\end{theorem}

\begin{theorem}[{cf. Kra et al., 2026, Theorem
2.5}]\protect\hypertarget{thm-FCP2}{}\label{thm-FCP2}

Let \(\AmenableGroup\) be an amenable group,
\(A\subset \AmenableGroup\), and \(\Folner\) be a Følner sequence in
\(\AmenableGroup\) such that the limit
\[\delta=\lim_{N\rightarrow\infty}\frac{\CountingMeasure{A\cap\Folner[N]}}{\CountingMeasure{\Folner[N]}}\]
exists.

Then there exists an ergodic measure-preserving system
\((\ShiftSpace,\Measure,\AmenableGroup)\) that is acted on by
\(\AmenableGroup\), a clopen set \(E\subset\ShiftSpace\), a Følner
sequence \(\FurstenbergFolner\) in \(\AmenableGroup\), and a point
\(a\in\ShiftSpace\) such that:

\begin{enumerate}
\def\labelenumi{\arabic{enumi}.}
\tightlist
\item
  The point \(a\) is generic with \(\Measure\) with respect to
  \(\FurstenbergFolner\).
\item
  \(\Measure(E)\geq\delta>0\).
\item
  \(A=\{\AmenableGroupElement\in \AmenableGroup:\GroupAction{\AmenableGroupElement}{a}\in E\}\).
\item
  Every measurable eigenfucntion is equal to \(\mu\)-almost everywhere
  to a topological eigenfunction.
\end{enumerate}

\end{theorem}

See:

(\citeproc{ref-kra2022infinite}{Kra et al., 2022}, Proposition 3.20 and
Lemma 5.7):

\begin{itemize}
\tightlist
\item
  Suppose for each \(j\in\mathbb{N}\), we have a system
  \((X_j,\Measure_j,T_j)\) and suppose there are continuous factor maps
  \(\pi_j:X_j\rightarrow X_{j-1}\) for all \(j>1\). Then there exists a
  unique \emph{inverse limit} \((X,\Measure,T)\), which is a system
  admitting continuous factor maps \(\psi_k:X\rightarrow X_j\) such that
  \(\pi_j\circ\psi_j=\psi_{j-1}\) for every \(j>1\), and such that
  \[\bigcup_{j\in\mathbb{N}}C(X_j)\circ\psi_j \] is dense in \(C(X)\).
  Additionally, whenever each of the systems \((X_j,\Measure_j,T_j)\) is
  ergodic, so is the inverse limit \((X,\Measure,T)\).

  \begin{itemize}
  \tightlist
  \item
    According to the original source {[}Nilpotent Structures in Ergodic
    Theory{]}, this applies to topological dynamical systems generally.
  \end{itemize}
\item
  Theorem 5.3: Fix an ergodic system \((X,\Measure,T)\). For every
  \(k\in\mathbb{N}\), there is a system \((Z_k,m_k,T)\) with all the
  following properties:

  \begin{itemize}
  \tightlist
  \item
    The system \((Z_k,m_k,T)\) is an inverse limit of ergodic \(k\)-step
    nilsystems.
  \item
    There is a measurable factor map \(\rho_k\) from \((X,\mu,T)\) to
    \((Z_k,m_k,T)\).
  \item
    For any \(f\in\text{L}^\infty(X,\Measure)\), we have
    \(||f||_{k+1}=0\) if and only if \(f\) and
    \(\text{L}^\infty(Z_k,m_k)\circ\rho_k\) are orthogonal in
    \(\text{L}^2(X,\Measure)\).
  \end{itemize}
\item
  Definition 5.4 \& 5.5:

  \begin{itemize}
  \tightlist
  \item
    Whenever \((X,\Measure,T)\) is ergodic, we call the system
    \((Z_k,m_k,T)\) satisfying the conclusion of Theorem 5.3 the
    \(k\)\emph{-step pronilfactor} of \((X,\Measure,T)\).
  \item
    Given an ergodic system \((X,\Measure,T)\) and \(k\in\mathbb{N}\),
    let \((Z_k,m_k,T)\) be its \(k\)-step pronilfactor. We say that
    \((X,\Measure,T)\) \emph{has topological pronilfactors} if, for
    every \(k\in\mathbb{N}\), there is a continuous factor map
    \(\pi_k:X\rightarrow Z_k\).
  \end{itemize}
\item
  Proposition 5.6 is equivalent to Charamaras and Mountakis
  (\citeproc{ref-charamaras2024findingproductsetsclasses}{2024}, Lemma
  3.5)
\item
  Lemma 5.7 is a generalisation of Proposition 3.20.

  \begin{itemize}
  \tightlist
  \item
    Both start with an ergodic system.
  \item
    Instead of \(a\in\text{gen}(\Measure,\Folner)\), \(a\) is a
    transitive point (admits dense orbit).
  \item
    The systems shown to exist:

    \begin{itemize}
    \tightlist
    \item
      Both be ergodic systems with Folner sequences with a generic point
      and a continuous factor map such that \(a\) maps to this generic
      point.
    \item
      Have topological pronilfactors in Lemma 5.7, instead of a
      continuous factor map to its Kronecker factor.
    \end{itemize}
  \item
    Lemma 5.7 also states the original ergodic system and the existing
    system are measure theoretically isomorphic systems.

    \begin{itemize}
    \tightlist
    \item
      Proposition 3.20 states this within its proof
    \end{itemize}
  \end{itemize}
\item
  Proof of Lemma 5.7 vs Proposition 3.20:

  \begin{itemize}
  \tightlist
  \item
    Measurable Factor Maps: Existence of measurable factor map/s to its
    pronilfactors or Kronecker Factor.
  \item
    Product/Kronecker system:

    \begin{itemize}
    \tightlist
    \item
      Lemma 5.7 creates a product space by taking the countable product
      over \(k\in\mathbb{N}\) of \(k\)-step pronilfactors.
    \item
      Proposition 3.20 focuses solely on the \(1\)-step pronilfactor,
      i.e., Kronecker factor.
    \item
      The transformation on these spaces is the product transformation
      of all \(k\)-step pronilfactors and is distal on this product
      system.
    \item
      Using a previous proposition, or essentially Charamaras and
      Mountakis
      (\citeproc{ref-charamaras2024findingproductsetsclasses}{2024},
      Lemma 3.5), to determine the generic point in the product space or
      Kronecker space.
    \end{itemize}
  \item
    New system:

    \begin{itemize}
    \tightlist
    \item
      A new system is created by taking the product between the original
      system and the Product/Kronecker system, defining the generic
      point as the product of the generic points in each system and the
      transformation similarly.
    \item
      The measure is defined as the push forward of the original measure
      with the measurable factor maps to each coordinate in the new
      system.
    \item
      Lemma 5.7 defines the new space to be the orbit closure of the
      generic point in the constructed product space and defining the
      projection of the first coordinate onto the original system.
    \end{itemize}
  \item
    The continuous map from the new system to the original system is
    defined as the projection on the first coordinate.

    \begin{itemize}
    \tightlist
    \item
      Similarly, the projection on the remaining coordinates are
      continuous maps from the new system to the original system's
      pronilfactors.
    \end{itemize}
  \item
    Isomorphism: The maps \(x\mapsto(x,\rho(x))\) and \((x,z)\mapsto z\)
    are inverses of each other almost everywhere and so the two spaces
    are measure theoretically isomorphic.

    \begin{itemize}
    \tightlist
    \item
      Thus, the new system is ergodic and the \(k\)-step pronilfactors
      of each system are isomorphic to each other.
    \end{itemize}
  \item
    Conclusion: the continuous maps from the new system to the
    original's \(k\)-step pronilfactors and the new system's generic
    point maps to the generic points in the \(k\)-step pronilfactors
    define continuous factor maps from the new system to its \(k\)-step
    pronilfactors.
  \end{itemize}
\end{itemize}

(\citeproc{ref-charamaras2024findingproductsetsclasses}{Charamaras and
Mountakis, 2024}, Lemma 3.5 \& Proposition 3.7):

\begin{itemize}
\tightlist
\item
  Lemma 3.5: Covered in Kronecker factor notes and already generalised
  to amenable groups. It is also a generalisation of Kra et al.
  (\citeproc{ref-kra2022infinite}{2022}, Propositions 3.19 \& 5.6).
\item
  Proposition 3.7: Similar generalising of Kra et al.
  (\citeproc{ref-kra2022infinite}{2022}, Proposition 3.20)
\end{itemize}

(\citeproc{ref-kra2025densityfinitesumstheorem}{Kra et al., 2025},
Definition 4.3 \& Subsection 4.4):

\begin{itemize}
\tightlist
\item
  Definition 4.3 (Nilsystems, pronilsystems): Let \(\Group\) be an
  \(s\)-step nilpotent Lie group and let \(\Gamma\subseteq\Group\) be a
  discrete and cocompact subgroup. The compact manifold
  \(X=\Group/\Gamma\) is an \(s\)-step nilmanifold.

  \begin{itemize}
  \tightlist
  \item
    The group \(\Group\) acts on \(X\) by left translation, and there is
    a unique Borel probability measure \(\Measure\) on \(X\) that is
    invariant under this action, the \emph{Haar measure} on \(X\).
    Letting \(T:X\rightarrow X\) denote left translation by a fixed
    element of \(\Group\), the resulting measure preserving system
    \((X,\Measure,T)\) is called an \(s\)-step nilsystem (as usual
    endowed with the Borel \(\sigma\)-algebra).
  \item
    An inverse limit of \(s\)-step nilsystems is called an \(s\)-step
    pronilsystem.
  \end{itemize}
\end{itemize}

Notes:

\begin{itemize}
\tightlist
\item
  This result is FCP but skipping the shift space straight to the system
  that has a continuous factor map to its Kronecker factor.

  \begin{itemize}
  \tightlist
  \item
    This is where property 4 comes from.
  \end{itemize}
\end{itemize}

\section{Comparison with Functional Analysis
Approach}\label{comparison-with-functional-analysis-approach}

\subsection{Defining the Algebra}\label{defining-the-algebra}

(\citeproc{ref-kra2007uniformity}{Kra and Host, 2007, p. 242}) uses a
separable subalgebra, \(\mathcal{A}\), which is a unitary subalgebra of
\(\ell^\infty(\mathbb{Z})\) that is invariant under the shift and under
complex conjugation, closed in \(\ell^\infty(\mathbb{Z})\) and separable
for the uniform norm written \(||\ \cdot\ ||_\infty\). They have the
intent to mostly consider the case of the separable subalgebra
\(\mathcal{A}(\boldsymbol{a})\) spanned by a bounded sequence
\(\boldsymbol{a}=(a_n:n\in\mathbb{Z})\). They also write \(\sigma\) for
the shift on \(\ell^\infty(\mathbb{Z})\) where, for a sequence
\(\boldsymbol{a}=(a_n:n\in\mathbb{Z})\),
\(\sigma\boldsymbol{a}=(a_{n+1}:n\in\mathbb{Z})\) and
\(\overline{\boldsymbol{a}}\) to denote the conjugate sequence
\((\overline{a}_n:n\in\mathbb{Z})\).

This means that:

\begin{itemize}
\tightlist
\item
  We have the constant sequence
  \(\boldsymbol{1}=(1:n\in\mathbb{Z})\in\mathcal{A}\).
\item
  For \(f\in\mathcal{A}\), we have
  \(||f||_\infty=\sup_{n\in\mathbb{Z}}|f_n|\),
  \(\sigma f\in\mathcal{A}\) and \(\overline{f}\in\mathcal{A}\).
\item
  For \(f,g\in\mathcal{A}\)\$, we have \(f+g\in\mathcal{A}\) and
  \(f\cdot g\in\mathcal{A}\).
\item
  If \(\boldsymbol{f}_n\in\mathcal{A}\) for all \(n\in\mathbb{N}\) and
  \(\lim_{n\rightarrow\infty}||f_n-f||_\infty=0\), then
  \(f\in\mathcal{A}\).
\item
  (Countable dense subset) There exists a countable set
  \(\{f_n:n\in\mathbb{N}\}\subseteq\mathcal{A}\) such that, for every
  \(f\in\mathcal{A}\) and every \(\varepsilon>0\), there exists
  \(n\in\mathbb{N}\) suchh that \(||f_n-f||_\infty<\varepsilon\).
\item
  \(\mathcal{A}(\boldsymbol{a})\) is generated by applying any
  combination of the above to \(\boldsymbol{a}\) and then taking the
  closure.
\end{itemize}

With our approach, the subalgebra is a consequence of our construction.
For example, \(E\subset\AmenableGroup\) corresponds to a vector where
\[e(\AmenableGroupElement)\begin{cases}1&\AmenableGroupElement\in E,\\0&\AmenableGroupElement\notin E. \end{cases} \]
This gives us
\[\boldsymbol{e}=(e(\AmenableGroupElement):\AmenableGroupElement\in\AmenableGroup)\in\ell^\infty(\AmenableGroup),\]
and \(\mathcal{A}(\boldsymbol{e})\) is the subalgebra we work with,
where we define the shift using the group action
\(\GroupAction{\AmenableGroupElement_1}{e(\AmenableGroupElement_2)}=e(\GroupOperation{\AmenableGroupElement_1}{\AmenableGroupElement_2})\).
Thus, the approach is equivalent so far.

\section{Defining the Dynamical
System}\label{defining-the-dynamical-system}

(\citeproc{ref-kra2007uniformity}{Kra and Host, 2007, p. 242}) goes on
to let \(X'\) be the \emph{Gelfand spectrum} of \(\mathcal{A}\), stating
this means \(X\) consists of the set of unitary homomorphisms from
\(\mathcal{A}\) to the complex numbers.

This means that: - For all \(f,g\in\mathcal{A}\),
\(\lambda\in\mathbb{C}\) and \(\varphi\in X'\), we have
\(\varphi(f+g)=\varphi(f)+\varphi(g)\),
\(\varphi(f\cdot g)=\varphi(f)\cdot \varphi(g)\), and
\(\varphi(\lambda f)=\lambda\varphi(f)\). - For all \(\varphi\in X\), we
have \(\varphi(\boldsymbol{1})=1\).

They then let \(C(X)\) denote the algebra of continuous functions on
\(X\) and conclude that there exists an isometric isomorphism of
algebras \(\Phi:C(X)\rightarrow\mathcal{A}\). They also define, for
\(\boldsymbol{b}\in\mathcal{A}\), the function \(\Phi^{-1}(b)\) as
\emph{the function associated to} \(\boldsymbol{b}\).

With our approach, we instead use the shift space
\(\ShiftSpace=\{0,1\}^\AmenableGroup\) such that
\(\boldsymbol{e}\in\ShiftSpace\), which ignores the subalgebra
construction, and define a topological dynamical system.

\protect\phantomsection\label{refs}
\begin{CSLReferences}{1}{0}
\bibitem[\citeproctext]{ref-charamaras2024findingproductsetsclasses}
Charamaras, D. and Mountakis, A. (2024). \emph{Finding product sets in
some classes of amenable groups}.
https://doi.org/\url{https://doi.org/10.1017/fms.2024.155}.

\bibitem[\citeproctext]{ref-kra2022infinite}
Kra, B. et al. (2022). \emph{Infinite sumsets in sets with positive
density}. Available at: \url{https://arxiv.org/abs/2206.01786}

\bibitem[\citeproctext]{ref-kra2025densityfinitesumstheorem}
Kra, B. et al. (2025). \emph{The density finite sums theorem}. Available
at: \url{https://arxiv.org/abs/2504.06424}

\bibitem[\citeproctext]{ref-kra2026shortproof}
Kra, B. et al. (2026). \emph{A short proof of erdős's \(B+C\)
conjecture}.

\bibitem[\citeproctext]{ref-kra2007uniformity}
Kra, B. and Host, B. (2007). \emph{Uniformity seminorms on
\(\ell^\infty\) and applications}. Available at:
\url{https://arxiv.org/abs/0711.3637}

\end{CSLReferences}




\end{document}
